\begin{proposition}[crossing and multiple-fibre modules]
\hypertarget{tag-03IP}{}\label{wbII:prop:relative-local}
\StatementRecord{wbII:prop:relative-local}{proved}{03IP}{the definitions and constructions stated in the proof}

The symbols $R_2$ and $R_3$ denote copies of the same regular local ring,
with different maps from the one-dimensional base:
let $R_2=\C\{x,y,z\}$ with base coordinate $t=xy$, and let
$R_3=\C\{x,y,z\}$ with base coordinate $t=xyz$.  Then
\[
\begin{aligned}
 M_2&:=\Omega^1_{R_2/\C\{t\}}
 =\frac{R_2\,dx\oplus R_2\,dy\oplus R_2\,dz}
        {R_2\,(y\,dx+x\,dy)},\\[2pt]
 M_3&:=\Omega^1_{R_3/\C\{t\}}
 =\frac{R_3\,dx\oplus R_3\,dy\oplus R_3\,dz}
        {R_3\,(yz\,dx+xz\,dy+xy\,dz)}.
\end{aligned}
\tag{W.176}\hypertarget{tag-03E6}{}\label{wbII:eq:CN:6.1}
\]
Both modules are torsion-free, but neither is locally free at the origin.
They embed in free modules by
\[
\begin{aligned}
 M_2&\hookrightarrow R_2e\oplus R_2dz,
 &dx&\longmapsto xe,&dy&\longmapsto-ye,&dz&\longmapsto dz,\\
 M_3&\hookrightarrow R_3e_1\oplus R_3e_2,
 &dx&\longmapsto xe_1,&dy&\longmapsto ye_2,&dz&\longmapsto-z(e_1+e_2),
\end{aligned}
\tag{W.177}\hypertarget{tag-03E7}{}\label{wbII:eq:CN:6.2}
\]
with images
\[
 R_2dz\oplus (x,y)e,
 \qquad
 R_3xe_1+R_3ye_2+R_3z(e_1+e_2).
\tag{W.178}\hypertarget{tag-03E8}{}\label{wbII:eq:CN:6.3}
\]
If $R_m=\C\{s,y_1,y_2\}$ and the base coordinate is $t=s^m$ with
$m\ge2$, then
\[
\Omega^1_{R_m/\C\{t\}}
\simeq R_m\,dy_1\oplus R_m\,dy_2
       \oplus R_m/(s^{m-1})\,ds,
\tag{W.179}\hypertarget{tag-03E9}{}\label{wbII:eq:CN:6.4}
\]
and the final summand is exactly its torsion submodule.
\end{proposition}



\begin{proof}

\begin{enumerate}[label={$\langle1\rangle$\arabic*.},leftmargin=*]
\item Differentiating the three equations gives the presentations in \eqref{wbII:eq:CN:6.1}.  For
$M_2$, the first map in \eqref{wbII:eq:CN:6.2} sends $(a,b,c)$ to $(ax-by,c)$.  If this is
zero, then $c=0$ and $ax=by$.  Since $x$ and $y$ are relatively prime in the
UFD $R_2$, there is an $h$ with $b=xh$ and $a=yh$; hence its kernel is exactly
the displayed relation.  For $M_3$, the second map sends $(a,b,c)$ to
$(ax-cz,by-cz)$.  If it is zero, relative primeness first gives $x\mid c$,
then $y\mid(c/x)$; writing $c=xyh$ gives
$(a,b,c)=h(yz,xz,xy)$, again exactly the defining relation.  Thus both maps
are injective, proving torsion-freeness and \eqref{wbII:eq:CN:6.3}.

\item At the maximal ideal $(x,y,z)$ each defining relation has all coefficients in
that maximal ideal.  Consequently the quotient of either module by the
maximal ideal has dimension $3$ over $\C$, while its generic rank is $2$;
a locally free module of rank $2$ would need only $2$ generators.  This proves
non-local-freeness at the origin.  Finally, for $t=s^m$ one has
$dt=m s^{m-1}ds=0$ relatively.  The integer $m$ is a unit in $\C$, so the
decomposition \eqref{wbII:eq:CN:6.4} follows.  Its last summand is annihilated by
$s^{m-1}$, while the other two summands are free; after inverting $s$ it
vanishes.  It is therefore exactly the torsion submodule.

\end{enumerate}
\end{proof}